% !TeX encoding = UTF-8
\documentclass[12pt,a4paper]{article}
%\usepackage{russ}

\usepackage[utf8]{inputenc}
\usepackage[english,russian]{babel}
\usepackage[left=3cm,right=3cm,top=3cm,bottom=3cm]{geometry}

\usepackage{amsmath,amsfonts,amssymb}
\usepackage{amscd}
\usepackage{amsthm}
%\usepackage{geometry}
%\usepackage{fancyhdr}
%\parindent = 0em %ðàññòîÿíèå îò êðàÿ â êðàñíûõ ñòðîêàõ
\usepackage{color}
\usepackage{cmap}

\usepackage{underscore} %нижние подчеркивания вне формул

%для ссылок
\usepackage{xcolor}
\usepackage{hyperref}
\definecolor{linkcolor}{HTML}{799B03} % цвет ссылок
%\definecolor{urlcolor}{HTML}{799B03} % цвет гиперссылок
\hypersetup{pdfstartview=FitH,  linkcolor=linkcolor, urlcolor=blue
	, colorlinks=true}

\usepackage{ulem} %зачеркивание via \sout{}

% чтобы itemize работал без лишних пробелов
\usepackage{enumitem}
\setlist{nolistsep,
	% itemsep=0.3cm,
	parsep=4pt}

\usepackage{color}
\newcommand{\tr}[1]{\textcolor{red}{#1}}

\newcommand{\todo}[1]{\marginpar{\scriptsize \tr{#1}}}

\newcommand{\sep}{
	\begin{center}
		\line(1,0){300}
	\end{center}
}

\begin{document}

\pagestyle{empty}

\section*{Концепция книги}
\subsection*{Особенности мира:}

Почти вся материя мира является \textit{живой} и есть магия, которая называется \textit{мастерством}. Мастерство позволяет более или менее произвольные манипуляции с живым веществом. В той или иной степени мастерством обладают все вплоть до растений. Люди сильно отличаются по мастерству, сильнейшие в боевом отношении равносильны армии обычных людей.
Существуют несколько форм телепатии, но чем более сложная информация передается, тем сильнее искажаются разумы обоих участников. Т. е. это не передача, а синхронизация с плохо контролируемыми побочными эффектами.
Также люди частично наследуют память и навыки родителей.

\subsection*{История мира}
Первый апокалипсис. Еще в незапамятную, первобытную эпоху случился первый апокалипсис. С неба упали демоны -- мертвые существа, превращающие все вещество в мертвое. Вся жизнь тогда чуть не погибла.

Золотой век. Через долгое время произошла технологическая революция: люди научились делать артефакты -- предметы, позволяющие применять заклинания, вложенные в них создателем. Применение артефактов позволило ''растиражировать'' таланты исключительных мастеров на все человечество. Произошел огромный рост населения, благосостояния и продолжительности жизни. Угнетение из-за материальных благ ушло в прошлое, сложилось анархическое, гуманистическое общество, сфокусированное на развитии уникальных индивидуальных способностей. Тесные телепатические контакты не приветствуются.

Второй апокалипсис. Золотой век кончился вторым апокалипсисом. Человечество выжило, но опять оказалось на грани исчезновения, несмотря на то, что его мощь со времен первого апокалипсиса невероятно возросла. Большинство людей оказалось абсолютно бесполезно и не участвовало в битве, а боеспособная верхушка общества наоборот полегла почти в полном составе. Немногие выжившие магистры вскоре погибли в междоусобице и большое количество самых важных технологий оказалось недоступно для использования. Вдобавок началась эпидемия \textit{чумы}. 

Эпоха чумы. Чума -- медленно и мучительно убивающая болезнь. Высокое мастерство позволяет сопротивляться ей ценой все возрастающих усилий, но в конечном итоге умирают все. Поэтому у человечества перестало хватать сил и времени на поддержание старых социальных институтов. Люди больше стали полагаться на телепатию для передачи навыков между быстро сменяющимися поколениями, обеспечения доверия в условиях беззакония и быстрого согласования интересов в коллективах (с помощью телепатии невозможно врать). Спустя некоторое количество поколений люди организуются в \textit{единодушия} -- группы, внутри которых постоянно поддерживается телепатический \textit{консенсус} -- общая для всех система целей и эмоциональных оценок. Единодушия редко превышают 200 человек из-за того сложности поддержания консенсуса.

\subsection*{Сюжет}

Умирающий от старости целитель Авиценна (на самом деле Полимегист -- выживший магистр золотого века) уговаривает Ходжу стать его учеником и продолжить очень важное дело, суть которого пока не ясна. Ходжа -- исключительно могущественный и умный маг, а также пророк нового религиозного движения. Полимегист считает, что только он способен продолжить его дело. В ходе обучения раскрывается информация об устройстве и истории мира. В конце Полимегист предлагает Ходже объединить разумы. Так он хочет объяснить свои планы и мотивы и одновременно убедиться, что Ходжа согласен с его планом. Герои видят воспоминания и мировоззрение друг друга.

Твист 1: Чуму наколдовал Полимегист, чтобы ускорить эволюцию человеческих способностей для подготовки к следующему апокалипсису.

Твист 2: Религиозное учение Ходжи о растворении в Боге оказывается не только мистическим откровением, но и вполне рациональным проектом объединения человечества в единый сверхразум. Сама природа людей изменилась за многие поколения жизни в единодушиях и подобное объединение является их естественным устремлением.

Растворение вызывает у Полимегиста экзистенциальный ужас и кажется равносильным исчезновению человечества. Чтобы помешать растворению он убивает Ходжу и прекращает чуму. Судьба человечества неясна.

Твист 3 (в книге будут только намеки): Герои книги -- это эволюционирующие протоколы, управляющие наномашинами в гигантском супе из наномашин, в который превращена их планета в ходе неудачной попытки колонизации расой прародителей. Апокалипсисы -- попытки прародителей исправить проблему.


\newpage
Дальше разная вспомогательная информация

\section*{Основные персонажи}
\begin{itemize}
	\item Полимегист (в юности Гермесион) -- великий маг и ученый, создатель чумы. Величайший артефактор, мастер телепортации и делатель богов. Потерял во время второго апокалипсиса почти все тело кроме мозга, заменив его на искусственное.
	
	Говорит ясно и академично, представитель рационализма в книге.
	
	\item Ходжа -- великий маг и ученый. Создатель теории построения больших единодуший. Верит в то, что все человеческие существа -- части одной божественной души. А построение единодуший приближает их участников к богу и служит способом получения божественных откровений. После окончания ученичества у Полимегиста получает имя Арист.

	Говорит поэтично, использует метафоры, сравнения, притчи. Но при этом в его словах всегда есть четкий смысл. На самом деле он тоже предельно рационален, но считает, что чтобы слова были восприняты, они должны иметь эмоциональную силу. Для него цель общения -- не представить мысль  собеседнику для рассмотрения его разумом, а отправить сообщение в самую душу.

	\item Арат - младший родственник фараона Хор-Птаха. Часть секретов фараонской династии содержится в его наследственной памяти. Начальник столичной охраны (власть его сильно ограничена ангелами, по факту он начальник полицейских структур, укомплектованных обычными людьми). Искусный ментальный лицемер, двоемыслец и ментальный самохирург -- умение развитое им, чтобы при сканировании разума казаться абсолютно лояльным. Арат применил техники самохирургии не только для создания маски, но и к "ядру" своей личности, чтобы стать лучшим человеком и избавиться от внутренних конфликтов (долг vs эгоизм) (рассказать об этом через память Агнии).
	\item Агния -- возлюбленная Полимегиста, потомок Арата. Изменила себя с помощью ментальной хирургии, чтобы стать совместимой с богом медицины и стать лучшей версией себя. Став аватаром, приняла имя Асклепия. Много веков сохраняет свои чувства к Полимегисту неизменными (отчасти намеренно, отчасти из-за ментальной ригидности аватаров). Ключевая фигура в перевороте Полимегиста и создании чумы (да и в преодолении второго апокалипсиса тоже).
	\item Александр -- величайший боевой маг, полководец и друг Полимегиста. Благодаря нему Содружество победило второй апокалипсис. Потом был убит Полимегистом потому что по этическим причинам не мог допустить ускорение эволюции через чуму.
\end{itemize}

\section*{Эволюционные изменения человечества эпохи чумы}

Точка по единству/принадлежности -- отчасти похоже на чувство потерявшегося ребенка к родителям, отчасти на кризис самоидентификации (когда не можешь понять, кто ты и чего хочешь, потому что твоя точка отсчета -- твой социум).

Люди теперь вечные дети поскольку постоянно живут в тени превосходящей их сущности (их единодушия). Единодушие воспринимается как родительская фигура, но также и как высшее существо (как у Дюркгейма).

Анимистичность -- склонность одушевлять сущности (как способ думать о чем-то, не обязательно всерьез).

Холистическое / интуитивное / анимическое мышление. Стало мало времени на академическое образование, философию, рассуждения и юридические споры. Зато на всех обрушилось огромное количество заимствованного опыта. Поэтому у людей развилась способность холистически мыслить. Вместо того, чтобы оперировать фактами и умозаключениями, люди научились лучше формировать внутренние ментальные модели (Гештальты?) явлений, процессов, сущностей. Эти модели неявные, неформализуемые, невыразимые словами, но действенные. Можно сказать что изучаемый объект воспринимается как живой, потом человек знакомится с ним через опыт и научается \textit{ощущать}, как этот объект будет вести себя.

Это анимическое мышление дает способность понимать сложные сущность вроде единодуший, настоящих богов золотого века (продолжающих в молчании свою деятельность), чумы, самой жизни как системы или демонов как цивилизации (при наличии данных).

\section*{Учение Ходжи}

Есть три уровня/способа понимания. Мистический (самый простой и 'поверхностный'), Схоластический (Философский) и Инструментальный (прагматический, рациональный).

\subsection*{Мистический}
Бог не материален и трансцендентен.\\
Он вечен, неизменен и предсуществует миру.\\
Бог -- это слово, учение и закон\\
Бог -- это все мы\\
Все человеческие души -- части божественной души\\
Душа разумного существа -- это дверь в мир идей.

Мир -- это ад и стонущая утроба. Бог воплотится в мире в виде человеческих душ.
И он постепенно пробуждается (и наши духовные практики приближают этот момент). 
Но одновременно пробуждаются злые сущности (как-бы их назвать джинны/демоны) и даже Ложные боги, искаженные грехами. Они придут не во плоти, а овладев человеческими телами. И будет великая (последняя?) война. Война материальная, но в первую очередь духовная. Но в итоге победит Бог и воздаст по заслугам праведникам и грешникам (а также их наследникам и потомкам).

Второе оживление.

Малое пророчество: демоны сойдут на 


Максимы:
\begin{itemize}
	\item К богу ведут множество путей.
	\item Человеческого разума не достаточно, чтобы постичь Бога, но можно коснуться его своим сознанием во время коллективного мистического ритуала.
	
\end{itemize}

\subsection*{Схоластический}
Единое.

Мир -- это зеркало (похоже на Платоновскую пещеру с тенями)

? Идея -- это информация преломленная в сознании.

\textit{Самовоспроизводящаяся идея}. Аналогия с живыми существами (бабочкой). То такое "бабочка"? Можно сказать что это набор физических объектов. Можно сказать -- что это идеальное понятие/идея объединяющее эти объекты. А можно сказать, что это алгоритм, описывающий процесс порождения новых бабочек из предшествующих бабочек. Или же 'генетический рецепт' лежащий определяющий этот алгоритм. Все эти вещи логически определяют друг друга. Поэтому можно считать что ``бабочка'' -- это одновременно конкретный набор существ, понятие, процесс, алгоритм или ``генетический рецепт''. Но что если этот ``рецепт'' бабочки -- всего лишь информация. Что если эта информация конечна и постижима? Тогда эта информация -- самовоспроизводящаяся идея.

Бог -- это идея. Самовоспроизводящаяся, Самоактуализирующаяся/самореализующаяся, универсальная (сама ее сущность такова, что она не может иметь версий), ретроактивная, превосходящая.
У Бога есть аспекты: идея бога = слово = учение, душа (изначально фрагментированная), воплощение (реальное существо в будущем).

\subsection*{Инструментальный}
Ретроактивная сущность. Конечное описание, универсальность. Эквивалентность пониманий.\\

Требования к Богу в терминах аспекта Учения:
\begin{itemize}
	\item защита от ереси
	\item мстительность (воздаяние после пробуждения)
	\item всеобщность (привлекательный компромисс включающий всех людей)
	\item ассимиляция (всегда можно добавить к учению еще людей с увеличением общего ``счастья'')
	\item конкурентность (по отношению к другим возможным учениям).
	\item самокооперция (разные воплощения не могут конфликтовать ни в какой форме и ``сливаются'' вместо этого; в каком-то смысле это невозможность копирования -- созвучно с мистическим принципом ``Бог един'')
\end{itemize}

Воздаяние после пробуждения воплощенного Бога -- печальное, но логически неизбежное свойство. Иные Боги невозможны т. к. они проигрывают своим мстительным двойникам.

\subsection*{Мысли для дискуссии с Полимегистом}

отсылки: пари Паскаля, Пантеизм, история про злой ретроактивный ИИ в блоге Элиезера.

Бог определяется не только своей идеей (исходных кодом / учением), но и средой в которой этот код / учение функционирует. Поэтому итоговое воплощение Бога зависит от всех нас, поскольку мы -- среда, в которой оно формируется и существует.

Важность разделенности на отдельный личности. Неправильность представления о ``слиянии'' людей в однородный суперорганизм (истинный бог -- это не масса которая растворяет и переваривает в себе людей). Потому что без этой отдельности не получается построить Бога, достойного любви и служения. Принципиальная распределенность этой системы и специально продуманные правила ее функционирования (учение или протоколы по которым действуют ее участники) дает возможность придать ей стабильность и консистентность. Нет одной точки принятия решений. Никто не может отринуть фундаментальные моральные принципы пока остальные их придерживаются. Упрощая, бог -- это толпа, которая не может свернуть с намеченного пути, потому что все связаны друг с другом. Также распределенность порождает внутреннюю конкуренцию.

Тело подобно дому для души. Раньше люди жили в своих разделенных домах им могли только видеть и слышать друг друга из окон. Только наследственная память была исключением из этой изоляции. Таким образом семья была обособленной социальной единицей. Когда появились единодушия, люди сбились в племена, у которых появилось общее пространство внутри которого они могли свободно общаться, сотрудничать и решать конфликты. Подобно деревне или племени древних охотников собирателей. Я построил первый город для душ. Сообщество настолько большое, что уже не может функционировать за счет личных отношений, договоренностей и освященных веками обычаев. Благодаря сложной организации незнакомцы могут работают сообща на общее благо, не испытывая проблем с доверием, порядком, соблюдением интересов каждого.

\section*{Понятия}
\begin{itemize}
	\item жизнь -- мистическая сила приводящая в движение материю. Практически все в достижимом мире является живым в той или иной степени. Магия и функционирование всех существ -- проявления жизни.
	\item мастерство/магия -- искусство манипулирования {\it живой} материей, перезаписывая волю жизненной силы находящейся в материи.
	\item нежизнь -- демонический аналог жизни, не являющейся ей. Альтернативная жизнь.
	\item демоны -- непостижимые неживые существа из космоса. По видимому стремятся уничтожить все живое.
	\item соединение разумов, прикосновение душ, сомышление, соразвитие?, комедитация (раньше была ментальная синхронизация) -- практики позволяющие обмениваться мыслями и навыками. Однако чем глубже и объемнее взаимодействие, тем сильнее меняется разум обоих участников. В частности некоторые способности могут быть разрушены подобной перестройкой разума. В золотом веке соединение разумов считалось отвратительным. Высочайшей ценностью считались уникальные способности, которые таким образом можно только сломать, но не приобрести.
	
	\item Содружество -- рыхлое дипломатической объединение разных государств, со временем покрывшее весь мир.
	\item Метрополис -- столица Содружества.
	\item Магистериум -- Правительство Содружества, включает в себя парламент магистров и различные министерства.
	\item магистр -- достаточно сильный маг, признанный Магистериумом. На момент основания Содружества было около 300 магистров, во время второго апокалипсиса -- более двух тысяч.
	\item Артефакт --  предмет, заклятый сопротивляться разложению жизнью подобно живым существам и, как правило, наделенный полезными свойствами.
	\item бог (золотого века) -- огромный артефакт глубоко интегрированный с мозгом мага, называемого называют аватаром этого бога. Создается с помощью контролируемого саморазрастания бога под контролем аватара до необъятной сущности в пределе охватывающей своим присутствием весь мир. Немногие имеют достаточно глубокое чувство магии в домене бога, чтобы стать его ядром/аватаром. Психика аватара становится малоподвижной так как он вынужден поддерживать синхронизацию с богом и подавлять изменения в своей психике, которые ей мешают или перестраивать самого бога.
	\item Божественный домен. На поверхностном уровне это область магии в которой сила аватара становится божественной. Но правильнее сказать домен -- часть самой природы (жизни / логоса) подвластная аватару. Так же доменом можно назвать сам способ манипуляции природными сущностями, за счет которого достигается божественная сила. Боги не склоны сосуществовать в одном домене, победитель получает весь домен.
	\item ментальная хирургия -- манипуляции над разумом при помощи артефактов. По сравнению с соединением разумов позволяет достичь более радикальных изменений и избежать обратного воздействия от пациента на хирурга. С другой стороны -- груба и рискованна для пациента. Применяется верхушкой Та-Кемет для насильственного приведения подданных к лояльности. В Содружестве считается абсолютным злом. Хотя наиболее тонкие и мягкие техники ментальной хирургии используются для создания богов.
	\item эрзац -- вспомогательный объект нужный для телепортации. При телепортации материя не перемещается. Вместо этого телепортируемый объект воссоздается в точке назначения из его эрзаца (одновременно и разрушается в точке отправления).
\end{itemize}


\section*{Мысли, которые хочется подсветить}
\begin{itemize}
	\item Эволюции по Дарвину и по Ламарку могут сосуществовать. Дарвиновский механизм требует изменений в генофонде, т. е. (в самом агрессивном и ускоренном варианте) смертей.
	\item Полимегист думал, что ''бытие определяет сознание''. И как только он вернет человечество в прежние условия, снова вернется и золотой век. Но оказалось что ''нельзя дважды войти в одну реку''. И измененные эпохой чумы люди и будут стремятся к растворению и без чумы.
	\item Цивилизация золотого века во многом похожа на постиндустриальное общество. Но в отличие от него, она построена на неявном знании. Отсюда усиленный по сравнению с современностью культ индивидуализма.
	\item Сильной стороной Полимегиста является то, что он являясь великим магом он является также великим ученым. При этом в своем плане ускоренной эволюции он жертвует научным развитием цивилизации, считая его менее важным, чем развитие врожденных способностей. Он не верит, что оно может сильно изменить цивилизацию (своего рода средневековое мышление).
	\item Успех действий Полимегиста (переворот и создание чумы) в конечном итоге основан на ментальной самохирургии Асклепии.
	\item Полимегист думал, что религиозность Ходжи -- просто следствие вынужденной необразованности всего общества. Он думал, что получив лучшее образование золотого века, Ходжа отбросит мистицизм. Но в итоге Полимегист сам оказался побежденным, увидев глубокую продуманность мировоззрения Ходжи, в которое полученные знания органично интегрировались.
	\item Из-за повсеместного использования объединения разумов, понимание человеческой души стало очень важной сферой деятельности (нужно как-то назвать). То что раньше было лишь отраслью философии стало очень практичной областью знания. Достижения человечества в ней сравнимы с достижениями науки во времена золотого века. Однако этот прогресс был не заметен для Полимегиста накопленные знания были субъективными и неявными.
	\item Перед объединением разумов Полимегист боялся того, что вынужден будет убить Ариста так как он окажется врагом его плана. Оказалось, что Арист был согласен с чумой, но из его мыслей следовало, что чума угрожает высшим ценностям самого Полимегиста.
	\item Высшие ценности Ходжи и Полимегиста прямо противоположны. Романтический индивидуализм против ультимативного коллективизма.
	\item новое человечество -- вечные дети своих единодуший.
	\item[?] Параллель логоса из стоицизма и жизни.
\end{itemize}


\section*{Принципы написания}
\begin{itemize}
	\item Быть кратким, придумать как опустить неинтересные детали.
	\item Показывать мир, а не объяснять.
	\item Нужно суметь заранее заинтриговать читателя сюжетными тайнами.
\end{itemize}

\subsection*{Проблемы}
\begin{itemize}
	\item Как назвать мастерство/магию. Магия мне не нравится тем, что имеет оттенок редкости/исключительности. А у меня ей владеет каждая собака. Мастерство как-то в тексте читается странно. К тому же нужно породить много производных понятий. Думаю, что придется оставить магию или придумать свое слово.
	\item Как говорить о Дарвновской/Ламарковской эволюции. Хочется чтобы это выглядело органично для фэнтези сеттинга, но при этом являлось настоящим научпопом (эдьютейнментом).
	\item Отсутствие главного вопроса в начале истории. Интрига не создана.
	\item Боги. Добавить про роль в обществе. И может, как они ощущаются. Чувствует ли Агния боль всего мира? Есть ли у нее ощущение всемогущества?
	\item Ходжа не раскрыт
	\item Отсутствие развития героев. Финальные портреты героев уже определены, можно создать развитие и заодно раскрытие, углубив предысторию:
	\begin{itemize}
		\item надлом Полимегиста вторым апокалипсисом
		\item как Агния решалась, как ощутила трансформацию
		\item описания 'моральных' битв Ходжи, его поступки в духе подставления второй щеки.
	\end{itemize}
	\item Добавить экшена (без детальных описаний, но эмоционально, с объяснением сути ситуации)
\end{itemize}

\subsection*{Не проблемы}
особенности книги, которые я не считаю проблемами.

\begin{itemize}
	\item Открытый финал -- оставлю! Проблема, по какому пути направить цивилизацию, -- сложная. И мне не хочется показывать некий путь как правильный или неправильный. Хочу подчеркнуть в финале, что Полимегист думал, что рассчитал правильный путь, но катастрофически ошибся. Мне хочется оставить читателя в конце с ощущением неясности и неуверенности.
	\item Персонажи — «носители идей». Стоит добавить в них глубины и психологизма, но в целом эта черта останется. Я пишу не про людей, а именно про эти идеи.
\end{itemize}

\section*{TODO}
\begin{itemize}
	\item Арат -> Тимолеон
	\item Читать авраамические притчи/проповеди про изгнание демонов, подстановку второй щеки, обращение перевоспитание и искупление грешников (Для Ходжи).
	\item Продумать и вплести отсылки и аналогии для понятий:
	\begin{itemize}
		\item идея общества как (внутреннего психологического) бога (Дюркгейм). Мб добавить фрейдистское Оно (или как там его).
		\item джихад и умма
		\item коллективные мистические обряды (радения?)
	\end{itemize}
\end{itemize}

\end{document}


